\documentclass[11pt]{article}
\usepackage[a4paper,margin=25mm]{geometry}
\usepackage{amsmath,amssymb,amsthm,mathtools,booktabs,hyperref,xurl,microtype}
\setlength{\emergencystretch}{3em}
\hypersetup{hidelinks}
\newtheorem{theorem}{Theorem}
\newtheorem{lemma}[theorem]{Lemma}
\newtheorem{corollary}[theorem]{Corollary}
\newcommand{\Q}{\mathbb Q}\newcommand{\Z}{\mathbb Z}
\newcommand{\re}{\operatorname{re}}
\title{Sharing the physical numerator chart\\in the ordinary-abc quartic transfer}
\author{\(\Z\Q\) Diophantine Project}
\date{24 November 2026}
\begin{document}\maketitle

\begin{abstract}
After sharing a physical numerator point and compressing reciprocal
products, its whole point chart can also be shared. Two field quadrics
describe each private point relative to the common numerator: the former
scale equation follows from an exact polynomial identity. One common
reference quotient removes the need for private scaled reciprocal
abscissas. The resulting presentation has 656 rational coordinates and
575 quadratic residuals, down from the explicit 674/593 predecessor.
Every new rational zero lifts polynomially to that complete predecessor,
even at zero values of the distinguished inputs. Completeness uses the
previously proved common-numerator arithmetic construction. All its
retained arithmetic premises, ordinary rational abc, and ineffective
defining coefficient remain unchanged.
\end{abstract}

\section{Scope and statement}

We start with the complete odd-Pell table of~\cite{odd}, incorporating
the common physical numerator theorem of~\cite{common}. The latter
uses the ordinary-abc integer interface of~\cite{paired}, the translated
secants of~\cite{secant}, and the arithmetic assertions retained there.
In particular the elliptic group, comparison, trace-index, prime-pattern,
local preparation, real recurrence and pole-realization assertions are
still premises. The sufficient odd positive coefficient remains
ineffective. This note is an algebraic refinement of that conditional
transfer; it supplies no new unconditional undecidability theorem over
\(\Q\).

\begin{theorem}\label{thm:main}
Under exactly the premises of~\cite{common}, its simultaneous ordinary-abc
integer interface admits a presentation with
\[
                44k+22\quad\text{rational coordinates},\qquad
                39k+17\quad\text{quadratic residuals}.
\]
The coordinates include the \(k\) distinguished inputs. Every rational
zero has integral inputs, and every tuple of nonzero integers has a
rational witness. Applied to the thirteen certified inputs of the
odd-Pell source~\cite{odd}, this gives a complete table with 656 rational
unknowns and 575 quadratic residuals. Its instance projection agrees
with that of the complete 674-variable predecessor.
\end{theorem}

The source instance is a coefficient, not another unknown. The source's
canonical positive witnesses, including its new odd-Pell half-quotient,
satisfy the nonzero-input condition needed for completeness. Soundness
does not assume that condition. The externally cited Jones coding
theorem remains required; no universality assumption for the historically
emended source is added.

\section{The two-equation private point chart}

Put
\[
 K=\Q(\rho),\qquad \rho^2=2,\qquad \Delta=11-6\rho,
 \qquad E:y^2=x^3-\Delta x.
\]
Every field coordinate means two rational coordinates in the basis
\(1,\rho\); \(\re(a+b\rho)=a\), and a bar denotes conjugation.
Write \(f(Q)=1/x(Q)\) and \(S_Q=y(Q)/x(Q)\). For the common physical
numerator \(U\), retain fields \(S=S_U\) and \(F=f(U)\). Its curve
equation is
\begin{equation}\label{eq:C}
                    C:=FS^2+\Delta F^2-1=0.
\end{equation}
For a private point \(V\), retain \(t=S_V\), \(G=f(V)\), and the
translated-secant field \(h\). Set \(\pi=h+S+t\) and impose
\begin{equation}\label{eq:AB}
 A:=\pi F-Gh=0,\qquad
 B:=\Delta(F+G)-(S-t)h+t^2=0.
\end{equation}

\begin{lemma}\label{lem:chart}
Equations~\eqref{eq:C}--\eqref{eq:AB} recover the full unscaled
translated-secant chart. In particular both reciprocal abscissas are
nonzero, the two points are distinct, and its difference output is
\(h\pi=x(U-V)\). No nonzero ordinate, chord parameter, or difference
abscissa is assumed.
\end{lemma}
\begin{proof}
The following identities hold as polynomials over \(K\):
\begin{align}
 Gh(S-t)-\Delta FG-1
   &=C-FB-(S-t)A, \label{eq:scale}\\
 Gt^2+\Delta G^2-1-C
   &=(G-F)B-(S-t)A. \label{eq:privatecurve}
\end{align}
Thus setting \(d=Gh\) gives all four original field equations
\[
 d=\pi F,\quad d=Gh,\quad
 \Delta(F+G)=(S-t)h-t^2,\quad
 1=d(S-t)-\Delta FG.
\]
The curve identities give \(F,G\ne0\). If \(F=G\) and \(S=t\),
then \(A=0\) gives \(S=0\), and \(B=0\) gives \(2\Delta F=0\),
a contradiction. The geometric interpretation, including tangent and
two-torsion fibers, now follows from the original secant chart
\cite[Proposition 2.1]{secant}.

For completeness, the inverse formulas are
\[
 h=\frac{(S+t)F}{G-F}\quad(G\ne F),\qquad
 h=\frac{\Delta(F+G)+t^2}{S-t}\quad(S\ne t).
\]
The curve identities make them agree on the overlap, and the two opens
cover all allowed distinct points. This also shows directly why an
individual zero ordinate creates no missing case.
\end{proof}

The common cubic equation~\eqref{eq:C} is kept quadratic by a shared
field \(S_2=S^2\): impose \(S_2-S^2=0\) and
\(FS_2+\Delta F^2-1=0\). Because these fields are shared, their cost
is independent of the number of inputs.

\section{The complete quadratic substitution}

Let \(X_i\) be the distinguished input. The predecessor retains
\(X_{2,i}=X_i^2\), \(X_{4,i}=X_{2,i}^2\), and scalar chart units
\(c_i\epsilon_i=1\). Its shared reference fields are
\[
 R=S_H,\quad T=S_{H-R_0},\quad v=v_H,
 \qquad R_0=(1+\rho,2\rho).
\]
The other reference elements are affine expressions
\[
 u=(\rho-1)(v-(4-2\rho)-R),\qquad
 w=(\rho-1)(v-(4-2\rho)+T).
\]
The retained reference chart makes \(u,w,v\) nonzero. Previously shared
fields \(J=f(H)f(U)\) and \(K=J^2\) remain, including the two
component equations for \(K-J^2\). Introduce a common reference
quotient \(\sigma\), a common product \(S_R\), and impose
\begin{equation}\label{eq:common}
       uF=vJ,\qquad w\sigma=v,\qquad S_R=SR.
\end{equation}
Together with the two field equations involving \(S_2\), these are
ten rational quadratic residuals.

For input \(i\), retain private fields \(t_i,G_i,h_i\), a scalar
\(q_i\), and a scalar \(\omega_i\), with
\begin{align}
 A_i&=B_i=0, &
 q_i&=\re\bigl(h_i(h_i+S+t_i)\bigr), \label{eq:private}\\
 a_i^{\rm old}&=X_{2,i}q_i, &
 t_i^{\rm old}&=X_it_i, &
 \omega_i&=c_iX_{2,i}. \label{eq:scaleout}
\end{align}
Here \(a_i^{\rm old}\) is the retained half-trace coordinate and
\(t_i^{\rm old}\) is the retained input-scaled ordinate-ratio field.
The reciprocal compression of~\cite{common} retained one private field
\(L_i\). Its second product relation becomes
\begin{equation}\label{eq:reciprocal}
                    \sigma\overline{L_i}=\omega_iG_i.
\end{equation}
These are eleven rational quadratic residuals per block. Keeping
\(\omega_i\) scalar costs one rational coordinate; retaining a scaled
field \(c_iX_i^2G_i\) would instead cost two.

If \(P_i,Q_i\) are the predecessor's two comparison fields, replace
its first comparison equation by
\begin{equation}\label{eq:comparison}
        2(X_{2,i}S_2-2X_iS_R+R^2)-P_i-c_iQ_i=0.
\end{equation}
This still consists of two rational quadrics. Its second comparison
equation is unchanged, using the retained \(t_i^{\rm old}\). All
height, contact, detector, pole, power-coding and source rows not
explicitly listed here remain literally unchanged after reindexing.

\section{Polynomial soundness, including zero inputs}

\begin{lemma}\label{lem:lift}
Every rational zero of the new table extends to a zero of the complete
674-variable predecessor, with exactly the same distinguished inputs
and source instance. No division by any \(X_i\) is required.
\end{lemma}
\begin{proof}
For each deleted coordinate use
\begin{align*}
 s_i^{\rm old}&=X_iS,& F_i^{\rm old}&=X_{2,i}F,&
 G_i^{\rm old}&=\omega_iG_i,\\
 h_i^{\rm old}&=X_ih_i,&
 d_i^{\rm old}&=\omega_iX_iG_ih_i,&
 E^{\rm old}&=vJ.
\end{align*}
These are polynomial expressions. Equations~\eqref{eq:scaleout}
also give the required already-retained scaled point and half-trace
coordinates. Use the retained equations \(X_{2,i}=X_i^2\),
\(X_{4,i}=X_i^4\), and \(c_i\epsilon_i=1\).
The old pair first-product, second-product, and secant residuals become,
respectively, \(X_i^3A_i\), zero, and \(-X_i^2B_i\). The old scale
residual becomes
\[
 c_iX_i^4\bigl(C-FB_i-(S-t_i)A_i\bigr)=0
\]
by~\eqref{eq:scale}. The half-trace residual is
\(X_i^2\bigl(q_i-\re(h_i(h_i+S+t_i))\bigr)=0\).

The first reciprocal-product residual is \(X_i^2(uF-vJ)=0\).
For the second, the exact identity
\begin{align*}
 v\overline{L_i}-w\omega_iG_i
  &=w(\sigma\overline{L_i}-\omega_iG_i)
        -(w\sigma-v)\overline{L_i}
\end{align*}
uses only the new rows. Finally \(S_2=S^2\), \(S_R=SR\) and
\(X_{2,i}=X_i^2\) identify~\eqref{eq:comparison} with
\(2(s_i^{\rm old}-R)^2-P_i-c_iQ_i=0\). The removed definition of
\(E\) holds by construction. Every remaining predecessor row is
unchanged and contains no deleted coordinate.

In particular these identities still hold when \(X_i=0\). They
give soundness on every new rational fiber, regardless of whether
the new restrictions admit witnesses for every zero-input tuple.
\end{proof}

\section{Completeness, counts and certificates}

For a nonzero integer tuple, the common-numerator theorem
\cite[Lemma 3]{common} chooses one physical \(U\), one reference
\(H\), and private \(V_i=U-[X_i]R_0\), with all required comparison
prime, fixed-place, real, contact and pole conditions. That theorem
uses the full all-sign order congruences and arbitrarily accurate local
approximation. We do not add any extra simultaneous arithmetic
constraint. Choose \(S,F,t_i,G_i,h_i\) from these points by
Lemma~\ref{lem:chart}, set \(\sigma=v/w\), \(S_2=S^2\),
\(S_R=SR\), \(q_i=\re(h_i(h_i+S+t_i))\) and
\(\omega_i=c_iX_i^2\). The old reciprocal relation, divided by
the reference unit \(w\), gives~\eqref{eq:reciprocal}. All new
equations hold. The pole arguments and values used for completeness
are exactly those in~\cite{common}.

In the tuple presentation, replace ten old rational point coordinates
per input by eight new private coordinates. Remove the two coordinates
of \(E\), and add the ten coordinates of \(S,F,S_2,S_R,\sigma\).
Starting from \(46k+14\) coordinates this gives
\[
       (46k+14)-10k-2+8k+10=44k+22.
\]
Thirteen old rows per block are replaced by eleven; two comparison
rows per block are changed but not added. The two old \(E\)-definition
rows are replaced by ten common rows. Thus the row count is
\[
                   (41k+9)-2k-2+10=39k+17.
\]
For thirteen inputs the saving is eighteen coordinates and eighteen
rows. Starting from the odd-Pell 674/593 table gives 656/575. This
proves Theorem~\ref{thm:main}. Completeness reselects elliptic
witnesses; it is not a bijection on arbitrary predecessor auxiliary
tuples.

Every residual has total degree at most two in the unknowns, with
integer polynomial coefficients in the retained parameters. The sum
of their squares has degree exactly four: the unchanged input-square
residual has a nonzero quadratic part, and squares of rational leading
quadratic forms cannot cancel.

The deterministic builder and independent checker are supplied in
this directory. Run from the repository root:
\begin{center}\small
\texttt{python -B}\enspace\nolinkurl{undecidability/abc_point_rechart_v105/build_point_rechart.py}\\
\texttt{python -B}\enspace\nolinkurl{undecidability/abc_point_rechart_v105/verify_point_rechart.py}.
\end{center}
The independent checker imports no builder. It interprets 197 old
rows, verifies 179 new or changed rows and 396 unchanged rows, audits
every occurrence of all 132 removed coordinates, and checks the eight
universal field identities used in the lift. It also verifies degree,
coefficient domain, full coordinate support, counts and hashes. These
finite checks certify the algebraic presentation and its reconstruction,
not the retained arithmetic premises or a numerical defining parameter.

\small
\begin{thebibliography}{9}\raggedright
\setlength{\itemsep}{2pt}
\bibitem{odd} \(\Z\Q\) Diophantine Project,
\emph{An odd second-Pell quotient removes one inverse}, 24 November 2026,
\nolinkurl{undecidability/odd_pell_quotient_v104/odd_pell_quotient.tex};
complete predecessor
\nolinkurl{undecidability/odd_pell_quotient_v104/odd_quotient_abc_quadrics.json}.
\bibitem{common} Project, \emph{A common numerator compresses reciprocal
products in the ordinary-abc quartic transfer}, 24 November 2026,
\nolinkurl{undecidability/abc_common_numerator_v103/common_numerator.tex}.
\bibitem{paired} Project, \emph{Paired contact coordinates: an ordinary-abc
quartic in 56 variables}, 19 November 2026,
\nolinkurl{papers/quartic_existence/paired_contact56_20261119_v81.tex}.
The paired chart and Laurent packing derive from Gabe Conant's notes
with Astra 6.2, audited in v78.
\bibitem{secant} Project, \emph{A quartic in 68 variables from translated
secants}, 17 November 2026,
\nolinkurl{papers/quartic_existence/quartic_abc68_20261117_v65.tex}.
\end{thebibliography}
\end{document}
